\chapter{Experimental: logical relations}
\label{chap:experimental-logrel}

\section*{At a glance}
\begin{itemize}
\item \lead{Goal} Prove definitional inversion for \texttt{SExpr} -- sort injectivity, $\Pi$
injectivity, and no sort is a $\Pi$ -- via a finitary semantic model, without the Church--Rosser
route of Chapter~\ref{chap:experimental-reduction}'s abstract theory.
\item \lead{Main results} Theorem~\ref{thm:expa-strongsound-uniq} (semantic unique typing);
Theorem~\ref{thm:expa-forallE-inv} ($\Pi$ injectivity); Theorem~\ref{thm:expa-sort-forallE-inv}
(sort $\neq$ $\Pi$); Theorem~\ref{thm:expa-sort-inv} (sort injectivity).
\item \lead{Status} \stExperimental{} a research module, off the verified path.
\stSorry{} Theorem~\ref{thm:expa-lr-adequacy}, the fundamental theorem, has a literal \texttt{sorry}
in its \texttt{const} case. \stTainted{} everything from Theorem~\ref{thm:expa-le-interp-strongsound}
on, including the three injectivity theorems above.
\item \lead{Lean files} \texttt{Experimental/ShapeLogRel.lean},
\texttt{Experimental/ShapeLogRelAdequacy.lean} (live); \texttt{Experimental/LogRel.lean},
\texttt{Experimental/CoinductiveLogRel.lean} (abandoned attempts, below).
\item \lead{Thesis} \texttt{unique.tex} (definitional inversion, unique typing) and
\texttt{soundness.tex} (the tagged-type ZFC model this development replaces).
\item \lead{Where} $\Gamma\vdash M\equiv N:A$ is \texttt{SExpr} definitional equality,
$\to^{*}_{\mathrm{wh}}$ its weak-head reduction (\texttt{WHRedS}); shapes order by $\le$, join by
$\sqcup$; $m:a$ is the shape-level typing judgement; all declarations live under a
\texttt{[Params]} instance in namespace \texttt{Lean4Lean.SExpr}.
\end{itemize}

\section{The shape algebra}

\begin{definition}[The shape hierarchy]
\label{def:expa-shape}
\lean{Lean4Lean.SExpr.Shape0,Lean4Lean.SExpr.ShapeS,Lean4Lean.SExpr.Shape,Lean4Lean.SExpr.ShapeFun}
\leanok
\srcloc{Lean4Lean/Experimental/ShapeLogRel.lean}{41}
\thesisref{\texttt{soundness.tex} line 325: the tagged types $T_n\subseteq U_n$ with universe index $n$.}

\lead{In short} A stratified family of \hl{finite approximations} of terms and types, indexed by a
level.
\begin{itemize}
\item $\mathtt{Shape}\ 0 = \mathtt{Shape0}$: only $\bot$ and $\mathtt{sort}\ r$ ($r=\mathtt{false}$
for $\mathrm{Prop}$, $\mathtt{true}$ for $\mathrm{Type}$).
\item $\mathtt{Shape}\ (n{+}1) = \mathtt{ShapeS}(\mathtt{Shape}\ n)$: adds $\mathtt{forallE}\ a\ f$,
$\mathtt{lam}\ f$, $\mathtt{ctor}\ c\ l$, $\mathtt{indTy}$.
\item $\mathtt{ShapeFun}\ n = \mathtt{List}(\mathtt{Shape}\ n\times\mathtt{Shape}\ n)$: a finite table
tabulating a monotone function on shapes.
\end{itemize}
\end{definition}

\begin{definition}[Level-polymorphic and smart constructors]
\label{def:expa-shape-ctors}
\lean{Lean4Lean.SExpr.Shape.bot,Lean4Lean.SExpr.Shape.sort,Lean4Lean.SExpr.Shape.prop,Lean4Lean.SExpr.Shape.type,Lean4Lean.SExpr.ShapeFun.bot,Lean4Lean.SExpr.Shape.lam',Lean4Lean.SExpr.Shape.ctor',Lean4Lean.SExpr.Shape.trim,Lean4Lean.SExpr.IsStruct}
\leanok
\srcloc{Lean4Lean/Experimental/ShapeLogRel.lean}{59}
\uses{def:expa-shape}
\lead{In short} $\bot$ and $\mathtt{sort}\ r$ exist at every level; \hl{smart constructors} collapse
uninformative shapes to $\bot$.
\begin{itemize}
\item $\mathtt{prop} = \mathtt{sort}\ \mathtt{false}$, $\mathtt{type} = \mathtt{sort}\ \mathtt{true}$,
$\mathtt{ShapeFun.bot} = [(\bot,\bot)]$.
\item $\mathtt{lam}'\ f = \bot$ when every output of $f$ is $\le\bot$.
\item $\mathtt{ctor}'\ c\ l = \bot$ when $c$ is an eta-constructor ($\mathtt{IsStruct}\ c$) and every
field of $l$ is $\le\bot$.
\item $\mathtt{trim}$ (\srcloc{Lean4Lean/Experimental/ShapeLogRel.lean}{820}) applies both
recursively.
\end{itemize}
\end{definition}

\begin{definition}[Compatibility of shapes]
\label{def:expa-shape-compat}
\lean{Lean4Lean.SExpr.Shape.Compat,Lean4Lean.SExpr.ShapeFun.Compat,Lean4Lean.SExpr.Shape.Compat.comm,Lean4Lean.SExpr.Shape.Compat.symm,Lean4Lean.SExpr.ShapeFun.Compat.def}
\leanok
\srcloc{Lean4Lean/Experimental/ShapeLogRel.lean}{72}
\uses{def:expa-shape}
\lead{In short} A decidable relation $s\mathbin{\mathtt{Compat}}t$: \hl{do two shapes agree wherever
both say anything}?
\begin{itemize}
\item $\bot$ is compatible with everything; two sorts, iff same flag; matching heads, iff components
are.
\item Tables: $\mathtt{Compat}\ f\ f'$ iff every compatible pair of inputs forces compatible outputs.
\item Symmetric (\texttt{comm}, \texttt{symm}); equivalent to a common upper bound only once shapes
are well formed (\ref{def:expa-wshape-join}).
\end{itemize}
\end{definition}

\begin{definition}[The information order on shapes]
\label{def:expa-shape-le}
\lean{Lean4Lean.SExpr.Shape.ble,Lean4Lean.SExpr.ShapeFun.ble,Lean4Lean.SExpr.Shape.LE,Lean4Lean.SExpr.ShapeFun.LE,Lean4Lean.SExpr.Shape.LE.def,Lean4Lean.SExpr.ShapeFun.LE.def,Lean4Lean.SExpr.Shape.LE.trans}
\leanok
\srcloc{Lean4Lean/Experimental/ShapeLogRel.lean}{123}
\uses{def:expa-shape,def:expa-shape-compat}
\lead{In short} A decidable \hl{preorder} $s\le t$, ``$t$ carries at least as much information as
$s$''.
\begin{itemize}
\item $\bot$ is least; matching heads compare componentwise.
\item Tables: $\mathtt{ble}\ f\ f'$ holds when every entry $(x,y)$ of $f$ is dominated by some entry
$(x',y')$ of $f'$ with $x'\le x$ and $y\le y'$.
\item \texttt{LE} and \texttt{DecidableRel} instances; reflexive and transitive.
\end{itemize}
\end{definition}

\begin{lemma}[Inversion and monotonicity for the shape order]
\label{family:expa-shape-le-lemmas}
\lean{Lean4Lean.SExpr.Shape.bot_le,Lean4Lean.SExpr.Shape.le_bot,Lean4Lean.SExpr.Shape.le_sort,Lean4Lean.SExpr.Shape.sort_le,Lean4Lean.SExpr.Shape.le_indTy,Lean4Lean.SExpr.Shape.indTy_le,Lean4Lean.SExpr.Shape.forallE_le,Lean4Lean.SExpr.Shape.forallE_le_forallE,Lean4Lean.SExpr.Shape.lam_le,Lean4Lean.SExpr.Shape.lam_le_lam,Lean4Lean.SExpr.Shape.LE.rfl,Lean4Lean.SExpr.Shape.LE.of_eq,Lean4Lean.SExpr.ShapeFun.LE.rfl,Lean4Lean.SExpr.ShapeFun.LE.trans,Lean4Lean.SExpr.Shape.Compat.mono,Lean4Lean.SExpr.Shape.Compat.mono_r,Lean4Lean.SExpr.ShapeFun.Compat.mono_r}
\leanok
\srcloc{Lean4Lean/Experimental/ShapeLogRel.lean}{143}

\lead{In short} Routine inversions: which shapes lie below or above one of a given head form.
\begin{center}
\begin{tabular}{p{0.43\linewidth}p{0.49\linewidth}}
\toprule
Statement & Reads \\
\midrule
$s\le\bot$ & iff $s=\bot$ \\
$s\le\mathtt{sort}\ r$ & iff $s=\bot$ or $s=\mathtt{sort}\ r$ \\
$\mathtt{sort}\ r\le s$ & iff $s=\mathtt{sort}\ r$ \\
$s\le\mathtt{indTy}$, $\mathtt{indTy}\le s$ & analogous \\
$\mathtt{forallE}$, $\mathtt{lam}$ pairs & analogous, componentwise \\
\bottomrule
\end{tabular}\par\smallskip
\end{center}
Plus reflexivity, transitivity, and monotonicity of \texttt{Compat} in the order.
\end{lemma}
\begin{proof}
\uses{def:expa-shape-le}
\leanok
\lead{Idea} \texttt{cases}/\texttt{simp} on the shape constructors, level $0$ handled separately from
$n{+}1$.
\end{proof}

\begin{definition}[Lifting a shape to another level]
\label{def:expa-shape-lift}
\lean{Lean4Lean.SExpr.Shape.lift,Lean4Lean.SExpr.ShapeFun.lift,Lean4Lean.SExpr.Shape.lift_self,Lean4Lean.SExpr.Shape.lift_lift,Lean4Lean.SExpr.Shape.lift_le_lift,Lean4Lean.SExpr.Shape.lift_mono,Lean4Lean.SExpr.Shape.Compat.lift,Lean4Lean.SExpr.Shape.lift_inj}
\leanok
\srcloc{Lean4Lean/Experimental/ShapeLogRel.lean}{284}
\uses{def:expa-shape,def:expa-shape-le}
\thesisref{\texttt{soundness.tex} line 325 ff.: the $(\mathtt{ulift},m,A)$ clause, lifting a tagged type without changing its meaning.}

\lead{In short} $\mathtt{lift}\ m$ rebuilds a shape at level $m$, \hl{truncating to $\bot$} what does
not fit.
\begin{itemize}
\item For $n\le m$: a full embedding -- identity at $n=m$, composes, reflects and preserves $\le$,
injective, preserves compatibility.
\end{itemize}
\end{definition}

\begin{definition}[Partial lowering of shapes]
\label{def:expa-shape-plift}
\lean{Lean4Lean.SExpr.Shape.plift,Lean4Lean.SExpr.ShapeFun.plift,Lean4Lean.SExpr.Shape.olift,Lean4Lean.SExpr.ShapeFun.olift,Lean4Lean.SExpr.Shape.plift_eq_lift,Lean4Lean.SExpr.Shape.plift_thm,Lean4Lean.SExpr.Shape.le_plift,Lean4Lean.SExpr.Shape.plift_le,Lean4Lean.SExpr.Shape.olift_thm}
\leanok
\srcloc{Lean4Lean/Experimental/ShapeLogRel.lean}{421}
\uses{def:expa-shape-lift}
\lead{In short} A \hl{partial inverse} of \texttt{lift}: the best approximation from below, plus an
exact preimage when one exists.
\begin{itemize}
\item $\mathtt{plift} : \mathtt{Shape}\ n \to \mathtt{Shape}\ m \times \mathtt{Option}(\mathtt{Shape}\
m)$; $\mathtt{olift}$ is its \texttt{Option} component.
\item Galois connection, for $n\le m$: $t \le (\mathtt{plift}\ s).1 \iff \mathtt{lift}_m\ t \le s$
(\texttt{le\_plift}); dually for \texttt{plift\_le} via the \texttt{Option} witness.
\item \texttt{set\_option backward.do.legacy true}, \srcloc{Lean4Lean/Experimental/ShapeLogRel.lean}{10}:
a Lean 4.32 \texttt{do}-elaborator compatibility switch these proofs need.
\end{itemize}
\end{definition}

\begin{definition}[Join of shapes]
\label{def:expa-shape-join}
\lean{Lean4Lean.SExpr.Shape.join,Lean4Lean.SExpr.ShapeFun.join,Lean4Lean.SExpr.ShapeFun.mem_join,Lean4Lean.SExpr.Shape.lift_join,Lean4Lean.SExpr.Shape.bot_join,Lean4Lean.SExpr.Shape.join_bot,Lean4Lean.SExpr.Shape.sort_join_sort}
\leanok
\srcloc{Lean4Lean/Experimental/ShapeLogRel.lean}{745}
\uses{def:expa-shape-le,def:expa-shape-compat}
\lead{In short} A total binary operation, \hl{well behaved only on compatible arguments}.
\begin{itemize}
\item $\bot$ a unit; two sorts join to themselves if flags agree, else $\bot$; matching heads join
componentwise; mismatched heads give $\bot$.
\item Tables: join collects, for every compatible input pair, the entry $(j, f(j)\sqcup f'(j))$ at
their joined input $j$ (\texttt{mem\_join} characterises membership); commutes with \texttt{lift}.
\item A least upper bound only on compatible arguments; proved at the well-formed level
(\ref{thm:expa-wshape-join-prop}).
\end{itemize}
\end{definition}

\begin{definition}[Well-formed shapes]
\label{def:expa-shape-wf}
\lean{Lean4Lean.SExpr.Shape.WF,Lean4Lean.SExpr.ShapeFun.WF,Lean4Lean.SExpr.ShapeFun.NonZero,Lean4Lean.SExpr.Shape.ListNonZero,Lean4Lean.SExpr.Shape.WF.lift_iff,Lean4Lean.SExpr.Shape.WF.olift}
\leanok
\srcloc{Lean4Lean/Experimental/ShapeLogRel.lean}{829}
\uses{def:expa-shape,def:expa-shape-join}
\lead{In short} The side conditions that make $\mathtt{lam}'$/$\mathtt{ctor}'$ collapse to $\bot$ --
\hl{how proof irrelevance enters the model}.
\begin{itemize}
\item A table is well formed when it: has an entry at input $\bot$; is closed under joins of
compatible inputs up to order-equivalence; is monotone; has well-formed components.
\item A shape is well formed when its tables are, plus: $\mathtt{lam}\ f$ needs $f$ \texttt{NonZero}
(some output not $\le\bot$); $\mathtt{ctor}\ c\ l$ needs $l$ \texttt{ListNonZero} when
$\mathtt{IsStruct}\ c$.
\item Preserved and reflected by \texttt{lift}.
\end{itemize}
\end{definition}

\begin{definition}[Well-formed shapes as a subtype, with a custom eliminator]
\label{def:expa-wshape}
\lean{Lean4Lean.SExpr.WShape,Lean4Lean.SExpr.WShapeFun,Lean4Lean.SExpr.WShape.bot,Lean4Lean.SExpr.WShape.sort,Lean4Lean.SExpr.WShape.indTy,Lean4Lean.SExpr.WShape.forallE,Lean4Lean.SExpr.WShape.lam,Lean4Lean.SExpr.WShape.lam',Lean4Lean.SExpr.WShape.ctor,Lean4Lean.SExpr.WShape.ctor',Lean4Lean.SExpr.WShapeFun.elems,Lean4Lean.SExpr.WShape.casesOn,Lean4Lean.SExpr.WShape.casesOn'}
\leanok
\srcloc{Lean4Lean/Experimental/ShapeLogRel.lean}{928}
\uses{def:expa-shape-wf}
\lead{In short} $\mathtt{WShape}\ n$, the subtype of well-formed shapes, with hand-written case
analysers that \hl{unpack the side conditions}.
\begin{itemize}
\item \texttt{Membership} instance, extensionality, lifted constructors ($\mathtt{lam}$ takes the
\texttt{NonZero} proof, $\mathtt{lam}'$ falls back to $\bot$).
\item $\mathtt{casesOn}'$ (level $n{+}1$) and $\mathtt{casesOn}$ (any level) present a well-formed
shape as $\bot$/sort/forallE/lam/ctor/indTy.
\item Restates the \texttt{Shape} order/lift/join API for the subtype
(\srcloc{Lean4Lean/Experimental/ShapeLogRel.lean}{1055}).
\end{itemize}
\end{definition}

\begin{definition}[Application of a shape function to a shape]
\label{def:expa-shape-app}
\lean{Lean4Lean.SExpr.ShapeFun.app,Lean4Lean.SExpr.Shape.app,Lean4Lean.SExpr.WShapeFun.app,Lean4Lean.SExpr.WShape.app,Lean4Lean.SExpr.ShapeFun.maxBelow,Lean4Lean.SExpr.ShapeFun.trunc,Lean4Lean.SExpr.WShapeFun.app_mono_l,Lean4Lean.SExpr.WShapeFun.app_mono_r,Lean4Lean.SExpr.WShape.app_mono_l,Lean4Lean.SExpr.WShape.app_mono_r,Lean4Lean.SExpr.TShape.app_mono,Lean4Lean.SExpr.WShape.Compat.app,Lean4Lean.SExpr.WShape.lam_app,Lean4Lean.SExpr.WShape.lam'_app}
\leanok
\srcloc{Lean4Lean/Experimental/ShapeLogRel.lean}{717}
\uses{def:expa-shape-join,def:expa-shape-le,def:expa-wshape}
\lead{In short} $f\ @\ x$: \hl{truncate, then take the maximum}.
\begin{itemize}
\item \texttt{trunc} keeps entries whose input is $\le x$; \texttt{maxBelow} returns the output of the
entry dominating all others, defaulting to $\bot$; for a well-formed table this maximum exists
(\texttt{ShapeFun.WF.app}, \srcloc{Lean4Lean/Experimental/ShapeLogRel.lean}{2270}).
\item Extended to \texttt{Shape} ($\bot$ unless the head is $\mathtt{lam}$), \texttt{WShape} and
\texttt{TShape}; monotone in both arguments, compatible with lifting and with joins.
\end{itemize}
\end{definition}

\begin{theorem}[Compatible well-formed shapes have a least upper bound]
\label{thm:expa-wshape-join-prop}
\lean{Lean4Lean.SExpr.WShape.join_prop,Lean4Lean.SExpr.WShape.join_prop.exists_max,Lean4Lean.SExpr.WShape.join_prop.wf_trunc,Lean4Lean.SExpr.WShape.join_prop.trunc_compat,Lean4Lean.SExpr.WShape.join_prop.app_core,Lean4Lean.SExpr.WShape.join_prop.of_compat,Lean4Lean.SExpr.WShape.join_prop.join_mem',Lean4Lean.SExpr.WShape.join_prop.compat_app_r,Lean4Lean.SExpr.WShape.join_prop.compat_app_l,Lean4Lean.SExpr.WShape.join_prop.ih_fun}
\leanok
\srcloc{Lean4Lean/Experimental/ShapeLogRel.lean}{1557}
\uses{def:expa-shape-join,def:expa-wshape}
\lead{In short} For $x,y:\mathtt{WShape}\ n$, \hl{compatibility and a common upper bound coincide},
and the join realises it.
\begin{itemize}
\item \lead{Given} $x$ and $y$ have a common upper bound.
\item \lead{Then} $x$ and $y$ are compatible; conversely, if compatible, $x.1\sqcup y.1$ is well
formed and $x.1\sqcup y.1\le z \iff x\le z\wedge y\le z$ for every well-formed $z$.
\end{itemize}
\end{theorem}
\begin{proof}
\uses{def:expa-shape-wf,def:expa-shape-app}
\leanok
\lead{Idea} Induction on the level. \lead{Steps}
\begin{enumerate}
\item Successor case factored into a local namespace
(\srcloc{Lean4Lean/Experimental/ShapeLogRel.lean}{1391}) with the inductive hypothesis as a section
variable.
\item Function-table case via a maximum-element argument and truncation.
\end{enumerate}
\end{proof}

\begin{definition}[Join on well-formed shapes and the \texttt{Join} relation]
\label{def:expa-wshape-join}
\lean{Lean4Lean.SExpr.WShape.join,Lean4Lean.SExpr.WShape.Join,Lean4Lean.SExpr.WShape.Compat.iff,Lean4Lean.SExpr.WShape.Join.mk,Lean4Lean.SExpr.WShape.Join.iff,Lean4Lean.SExpr.WShapeFun.join,Lean4Lean.SExpr.WShapeFun.Join,Lean4Lean.SExpr.WShapeFun.ofElems}
\leanok
\srcloc{Lean4Lean/Experimental/ShapeLogRel.lean}{1627}
\uses{thm:expa-wshape-join-prop}
\lead{In short} $x\sqcup y$ on well-formed shapes: the underlying join if compatible, else $\bot$ --
\hl{always well formed}.
\begin{itemize}
\item Relational form $\mathtt{Join}\ x\ y\ z$: $\forall w,\ z\le w\iff x\le w\wedge y\le w$.
\item \texttt{Join.mk} builds it from compatibility; \texttt{Join.iff} translates between the two
forms; \texttt{Compat.iff} reads compatibility off a common upper bound.
\end{itemize}
\end{definition}

\begin{definition}[Level-erased shapes]
\label{def:expa-tshape}
\lean{Lean4Lean.SExpr.TShape,Lean4Lean.SExpr.WShape.T,Lean4Lean.SExpr.TShape.LE,Lean4Lean.SExpr.TShapeFun.LE,Lean4Lean.SExpr.TShape.LE.def,Lean4Lean.SExpr.TShape.bot,Lean4Lean.SExpr.TShape.sort,Lean4Lean.SExpr.TShape.Compat,Lean4Lean.SExpr.TShape.join,Lean4Lean.SExpr.TShape.Join,Lean4Lean.SExpr.TShape.LE.forallE_decomp,Lean4Lean.SExpr.TShape.LE.lam'_decomp,Lean4Lean.SExpr.TShape.LE.le_lam}
\leanok
\srcloc{Lean4Lean/Experimental/ShapeLogRel.lean}{1819}
\uses{def:expa-wshape,def:expa-shape-lift,def:expa-wshape-join}
\thesisref{\texttt{soundness.tex} line 325 ff.: $\mathtt{TShape}$ is $\bigcup_n T_n$, its order mirroring the $(\mathtt{ulift},m,A)$ clause.}

\lead{In short} $\mathtt{TShape}=\Sigma n,\ \mathtt{WShape}\ n$ \hl{forgets the level}.
\begin{itemize}
\item Order, compatibility, join: lift both arguments to the max of their levels; any common upper
bound of the levels works instead (\texttt{LE.def}).
\item Decomposition lemmas turn $\le$ between two $\mathtt{forallE}$ (resp.\ $\mathtt{lam}'$) shapes
into componentwise inequalities after lifting.
\item $\mathtt{TShape.type}$ exists (\srcloc{Lean4Lean/Experimental/ShapeLogRel.lean}{1887}); no
$\mathtt{TShape.prop}$ counterpart.
\end{itemize}
\end{definition}

\begin{lemma}[Head-form discrimination for level-erased shapes]
\label{family:expa-tshape-discrimination}
\lean{Lean4Lean.SExpr.TShape.sort_not_le_lam',Lean4Lean.SExpr.TShape.forallE_not_le_lam',Lean4Lean.SExpr.TShape.ctor_not_le_lam',Lean4Lean.SExpr.TShape.indTy_not_le_lam',Lean4Lean.SExpr.TShape.lam_not_le_forallE,Lean4Lean.SExpr.TShape.sort_not_le_forallE,Lean4Lean.SExpr.TShape.sort_not_le_ctor',Lean4Lean.SExpr.TShape.forallE_not_le_ctor',Lean4Lean.SExpr.TShape.lam_not_le_ctor',Lean4Lean.SExpr.TShape.indTy_not_le_forallE,Lean4Lean.SExpr.TShape.indTy_not_le_sort,Lean4Lean.SExpr.TShape.ctor_not_le_forallE}
\leanok
\srcloc{Lean4Lean/Experimental/ShapeLogRel.lean}{1928}

\lead{In short} Shapes with \hl{distinct non-$\bot$ head constructors} are incomparable in the
\texttt{TShape} order: twelve instances, e.g.\ $\neg(\mathtt{sort}\ r)\le\mathtt{lam}'\ f$,
$\neg(\mathtt{forallE}\ a\ f_1)\le\mathtt{lam}'\ f_2$.
\begin{itemize}
\item \lead{Caveat} Not exhaustive: e.g.\ \texttt{indTy\_not\_le\_ctor'} and
\texttt{ctor\_not\_le\_sort} are missing; only the pairs later proofs need are covered.
\end{itemize}
\end{lemma}
\begin{proof}
\uses{def:expa-tshape}
\leanok
\lead{Idea} Unfold \texttt{TShape.LE.def} at a common level, rewrite the lifted head forms, and
\texttt{simp} \texttt{Shape.ble} to \texttt{False}; the $\mathtt{lam}'$ cases split on smart-constructor
collapse.
\end{proof}

\begin{definition}[Singleton shape functions]
\label{def:expa-shape-single}
\lean{Lean4Lean.SExpr.ShapeFun.single,Lean4Lean.SExpr.WShapeFun.single,Lean4Lean.SExpr.ShapeFun.WF.single,Lean4Lean.SExpr.WShapeFun.single_app,Lean4Lean.SExpr.WShapeFun.single_le,Lean4Lean.SExpr.WShapeFun.compat_single,Lean4Lean.SExpr.WShapeFun.Join.app_l}
\leanok
\srcloc{Lean4Lean/Experimental/ShapeLogRel.lean}{2405}
\uses{def:expa-shape-app,def:expa-shape-wf}
\lead{In short} $\mathtt{single}\ x\ y=[(x,y)]$ (plus $(\bot,\bot)$ unless $x=\bot$): the \hl{least
table sending $x$ to $y$}.
\begin{itemize}
\item Well formed; $(\mathtt{single}\ x\ y)\ @\ x' = y$ if $x\le x'$, else $\bot$; below any
well-formed table with that behaviour.
\item Used to build witnesses in the soundness proofs (\ref{thm:expa-le-interp-sound-cases}).
\end{itemize}
\end{definition}

\section{The shape-level typing judgement}

\begin{definition}[Shape typing]
\label{def:expa-shape-hastype}
\lean{Lean4Lean.SExpr.hasType.core,Lean4Lean.SExpr.Shape.hasType,Lean4Lean.SExpr.Shape.HasType,Lean4Lean.SExpr.Shape.HasDom,Lean4Lean.SExpr.Shape.HasTypePi,Lean4Lean.SExpr.Shape.HasTypeLam,Lean4Lean.SExpr.WShape.HasType,Lean4Lean.SExpr.WShape.HasDom,Lean4Lean.SExpr.WShape.HasTypePi,Lean4Lean.SExpr.WShape.HasTypeLam,Lean4Lean.SExpr.TShape.HasType,Lean4Lean.SExpr.TShape.HasTypeLam}
\leanok
\srcloc{Lean4Lean/Experimental/ShapeLogRel.lean}{2482}
\uses{def:expa-shape-app,def:expa-shape-le,def:expa-wshape}
\thesisref{\texttt{soundness.tex} line 325 ff.: the clause deriving $(\Pi,A,B)\in T_n$ from $A\in T_n$, $B:\mathrm{tag}(A)\to T_n$.}

\lead{In short} A decidable relation $m:a$, ``\hl{$m$ inhabits the type-shape $a$}'', by simultaneous
recursion; kept boolean only for decidability -- Theorem~\ref{thm:expa-hastype-unfold} tabulates the
rules used in proofs.
\begin{itemize}
\item \texttt{HasDom}, \texttt{HasTypePi}, \texttt{HasTypeLam} name the table-level side conditions;
lifted to \texttt{WShape} and, by lifting to a common level, to \texttt{TShape}.
\end{itemize}
\end{definition}

\begin{theorem}[Inductive characterisation of shape typing]
\label{thm:expa-hastype-unfold}
\lean{Lean4Lean.SExpr.Shape.HasTypeU,Lean4Lean.SExpr.Shape.HasType.unfold,Lean4Lean.SExpr.Shape.HasType.unfold_iff,Lean4Lean.SExpr.WShape.HasTypeU,Lean4Lean.SExpr.WShape.HasType.unfold,Lean4Lean.SExpr.WShape.HasType.unfold_iff}
\leanok
\srcloc{Lean4Lean/Experimental/ShapeLogRel.lean}{2513}
\uses{def:expa-shape-hastype}
\lead{In short} The boolean relation is equivalent to an inductive judgement \texttt{HasTypeU}, one
rule per head form -- \hl{the inversion principle used everywhere below}.
\begin{center}
\begin{tabular}{lp{0.79\linewidth}}
\toprule
Head form & Rule \\
\midrule
$\bot$ & $\bot:x$ from $x:\mathrm{Type}$ \\
$\mathtt{sort}\ r$ & $\mathtt{sort}\ r:\mathrm{Type}$ \\
$\mathtt{forallE}\ a\ b$ & from \texttt{HasTypePi}, at $\mathtt{sort}\ r$ \\
$\mathtt{lam}\ f$ & from \texttt{HasTypeLam}, at $\mathtt{forallE}\ a\ b$ \\
$\mathtt{ctor}\ c\ l$ & $:\mathtt{indTy}$ \\
$\mathtt{indTy}$ & $:\mathrm{Type}$ \\
\bottomrule
\end{tabular}\par\smallskip
\end{center}
\end{theorem}
\begin{proof}
\leanok
\lead{Idea} Forwards: \texttt{split} on the defining match, \texttt{constructor}, \texttt{simp}. Backwards:
case analysis on the derivation, the $\bot$ case recursing on its premise's inversion.
\end{proof}

\begin{theorem}[Monotonicity of shape typing]
\label{thm:expa-hastype-mono}
\lean{Lean4Lean.SExpr.WShape.HasType.mono_r,Lean4Lean.SExpr.WShape.HasType.mono_l,Lean4Lean.SExpr.WShape.HasDom.mono_r,Lean4Lean.SExpr.WShape.HasDom.mono_l,Lean4Lean.SExpr.WShape.HasTypeLam.mono_r,Lean4Lean.SExpr.WShape.HasTypePi.mono_l,Lean4Lean.SExpr.WShape.HasTypeLam.mono_l,Lean4Lean.SExpr.TShape.HasType.mono_r,Lean4Lean.SExpr.TShape.HasType.mono_l,Lean4Lean.SExpr.find_cycle,Lean4Lean.SExpr.WShape.HasType.mono_l.ih_dom,Lean4Lean.SExpr.WShape.HasType.mono_l.ih_pi,Lean4Lean.SExpr.WShape.HasType.mono_l.ih_lam}
\leanok
\srcloc{Lean4Lean/Experimental/ShapeLogRel.lean}{2681}
\uses{thm:expa-hastype-unfold,def:expa-shape-le}
\lead{In short} Two monotonicity principles, one needing \hl{order-equivalence, not just $\le$}.
\begin{itemize}
\item \lead{Given} $m:a$, $a\le a'$, $a':\mathtt{sort}\ r$. \lead{Then} $m:a'$ (\texttt{mono\_r}).
\item \lead{Given} $m:a$, $m\le m'$ \emph{and} $m'\le m$. \lead{Then} $m':a$ (\texttt{mono\_l}).
\item Both restated for \texttt{TShape}.
\end{itemize}
\end{theorem}
\begin{proof}
\uses{def:expa-shape-app}
\leanok
\lead{Idea} \texttt{mono\_r}: induction on the level via \texttt{HasTypeU} inversion. \texttt{mono\_l}: a
pigeonhole lemma \texttt{find\_cycle} (\srcloc{Lean4Lean/Experimental/ShapeLogRel.lean}{2716}, a
finite list under a transitive relation where every element points somewhere reaches a reflexive
point).
\end{proof}

\begin{theorem}[Shape typing is closed under joins]
\label{thm:expa-hastype-join}
\lean{Lean4Lean.SExpr.WShape.HasType.join,Lean4Lean.SExpr.WShape.HasDom.join,Lean4Lean.SExpr.WShape.HasTypePi.join,Lean4Lean.SExpr.WShape.HasTypeLam.join,Lean4Lean.SExpr.WShape.HasType.join',Lean4Lean.SExpr.TShape.HasType.join,Lean4Lean.SExpr.TShape.HasType.join'}
\leanok
\srcloc{Lean4Lean/Experimental/ShapeLogRel.lean}{2960}
\uses{thm:expa-hastype-unfold,def:expa-wshape-join}
\lead{In short} The shapes approximating a term form a \hl{directed set with suprema}.
\begin{itemize}
\item \lead{Given} $m_1:a$, $m_2:a$, $m_1$ compatible with $m_2$.
\item \lead{Then} $m_1\sqcup m_2:a$; primed variants take the join relationally, and the type-shape
may be joined too.
\end{itemize}
\end{theorem}
\begin{proof}
\uses{thm:expa-wshape-join-prop,def:expa-shape-app}
\leanok
\lead{Idea} Induction on the level, case analysis on head forms, using \ref{thm:expa-wshape-join-prop},
\texttt{HasTypeU} inversion and table-membership lemmas; domain and lambda cases factored into local
recursors.
\end{proof}

\begin{theorem}[Shape-level proof irrelevance and retyping]
\label{thm:expa-hastype-proofirrel}
\lean{Lean4Lean.SExpr.WShape.HasType.proofIrrel,Lean4Lean.SExpr.TShape.HasType.proofIrrel,Lean4Lean.SExpr.WShape.HasType.retype,Lean4Lean.SExpr.TShape.HasType.retype,Lean4Lean.SExpr.TShape.LE.le_forall,Lean4Lean.SExpr.LE_Forall,Lean4Lean.SExpr.TShape.HasType.ty_forallE_inv}
\leanok
\srcloc{Lean4Lean/Experimental/ShapeLogRel.lean}{3211}
\uses{thm:expa-hastype-unfold,def:expa-shape-wf}
\thesisref{\texttt{axioms.tex}, proof irrelevance; \texttt{unique.tex} \S\texttt{sec:kappa}, where it breaks Church--Rosser -- here it is absorbed into the model instead.}

\lead{In short} \hl{Every inhabitant of a $\mathrm{Prop}$-shape is $\bot$}.
\begin{itemize}
\item \lead{Given} $a:\mathtt{sort}\ \mathtt{false}$, $x:a$. \lead{Then} $x=\bot$
(\texttt{proofIrrel}).
\item \lead{Given} $a:\mathtt{sort}\ r$, $a':\mathtt{sort}\ r'$, $a\le a'$. \lead{Then}
$a:\mathtt{sort}\ r'$ (\texttt{retype}).
\item \texttt{TShape.LE.le\_forall} inverts $\le$ against $\mathtt{forallE}$ via the auxiliary
\texttt{LE\_Forall}; \texttt{ty\_forallE\_inv}: anything inhabiting a $\mathtt{forallE}$ shape is
$\bot$ or $\mathtt{lam}'$.
\end{itemize}
\end{theorem}
\begin{proof}
\leanok
\lead{Idea} \texttt{HasTypeU} inversion on both premises: level $0$ forces $\bot$; at level $n{+}1$ the
induction hypothesis forces every output of the lambda's table to $\bot$, so $\mathtt{lam}'$
collapses it.
\end{proof}

\section{The shape interpretation}

\begin{definition}[Shape valuations for de Bruijn contexts]
\label{def:expa-valuation}
\lean{Lean4Lean.SExpr.Valuation,Lean4Lean.SExpr.Valuation.nil,Lean4Lean.SExpr.Valuation.push,Lean4Lean.SExpr.Valuation.LE,Lean4Lean.SExpr.Valuation.Compat,Lean4Lean.SExpr.Valuation.join}
\leanok
\srcloc{Lean4Lean/Experimental/ShapeLogRel.lean}{3282}
\uses{def:expa-tshape}
\lead{In short} $\mathtt{Valuation}=\mathbb{N}\to\mathtt{TShape}$: \hl{a level-erased shape per de
Bruijn index}.
\begin{itemize}
\item \texttt{nil} sends everything to $\bot$; \texttt{push} extends at index $0$.
\item Order, compatibility, join: pointwise; \texttt{LE.push} splits the order on an extended
valuation.
\end{itemize}
\end{definition}

\begin{definition}[The shape interpretation relation]
\label{def:expa-le-interp}
\lean{Lean4Lean.SExpr.LE_Interp,Lean4Lean.SExpr.LE_Interp.Matches,Lean4Lean.SExpr.LE_Interp.RHS,Lean4Lean.SExpr.LE_Interp.Const}
\leanok
\srcloc{Lean4Lean/Experimental/ShapeLogRel.lean}{3333}
\uses{def:expa-valuation,def:expa-shape-hastype,def:expa-shape-app,class:expb-sexpr-params,inductive:pattern-core}
\thesisref{\texttt{soundness.tex} \S Modeling Lean in ZFC (line 118) and \S Type injectivity (line 325): a finitary, one-sided syntactic stand-in for $m\le[\![\Gamma\vdash M]\!]_\gamma$.}

\lead{In short} $\mathtt{LE\_Interp}\ \rho\ m\ M$: \hl{$m$ is a lower bound for the meaning of $M$}
under valuation $\rho$.
\begin{center}
\begin{tabular}{lp{0.84\linewidth}}
\toprule
Term $M$ & Rule \\
\midrule
any & $\bot$ interprets anything \\
$\mathtt{bvar}\ i$ & $m\le\rho\ i$ \\
$\mathtt{sort}\ \ell$ & $m\le\mathtt{sort}(\ell\neq 0)$ \\
$F\ A$ & $m\le f\ @\ a$, from shapes $f$ of $F$, $a$ of $A$ \\
$\lambda A.\,t$ & $m\le\mathtt{lam}'\ f$: domain $A$ has shape $a$, $f$ has domain $a$, each output
$f\ @\ x$ interprets the body under $\rho.\mathtt{push}\ x$ \\
$\forall A.\,B$ & $m\le\mathtt{forallE}\ b\ f$, analogous, with domain shapes $b,b'$ \\
constant & from \texttt{Const}/\texttt{Matches}/\texttt{RHS}, using the declared type and the
environment's $\delta$/$\iota$ pattern rules \\
\bottomrule
\end{tabular}\par\smallskip
\end{center}
\end{definition}

\begin{lemma}[Smart constructors and inversions for the interpretation]
\label{family:expa-le-interp-intro}
\lean{Lean4Lean.SExpr.LE_Interp.bvar',Lean4Lean.SExpr.LE_Interp.bvar0,Lean4Lean.SExpr.LE_Interp.sort',Lean4Lean.SExpr.LE_Interp.app',Lean4Lean.SExpr.LE_Interp.lam',Lean4Lean.SExpr.LE_Interp.forallE',Lean4Lean.SExpr.LE_Interp.bvar_iff,Lean4Lean.SExpr.LE_Interp.le_sort,Lean4Lean.SExpr.LE_Interp.le_sort'}
\leanok
\srcloc{Lean4Lean/Experimental/ShapeLogRel.lean}{3352}

\lead{In short} Specialisations at reflexivity, plus \hl{inversion at \texttt{bvar} and sort}.
\begin{itemize}
\item $\mathtt{LE\_Interp}\ \rho\ m\ (\mathtt{bvar}\ i)\iff m\le\rho\ i$.
\item At a sort: $m\le\mathtt{sort}(\ell\neq 0)$.
\end{itemize}
\end{lemma}
\begin{proof}
\uses{def:expa-le-interp}
\leanok
\lead{Idea} Smart constructors: immediate from the rules. Inversions: one-step case analysis, using that
the $\bot$ rule always yields a shape below $\bot$.
\end{proof}

\begin{theorem}[Monotonicity of the shape interpretation]
\label{thm:expa-le-interp-mono}
\lean{Lean4Lean.SExpr.LE_Interp.mono,Lean4Lean.SExpr.LE_Interp.mono_l,Lean4Lean.SExpr.LE_Interp.unlift,Lean4Lean.SExpr.LE_Interp.RHS.mono,Lean4Lean.SExpr.LE_Interp.Const.mono,Lean4Lean.SExpr.LE_Interp.RHS.mono_l,Lean4Lean.SExpr.LE_Interp.Matches.mono_l,Lean4Lean.SExpr.LE_Interp.Const.mono_l}
\leanok
\srcloc{Lean4Lean/Experimental/ShapeLogRel.lean}{3378}
\uses{def:expa-le-interp,def:expa-tshape}
\lead{In short} \hl{Downward closed} in the shape, monotone in the valuation.
\begin{itemize}
\item $m'\le m$ and $\mathtt{LE\_Interp}\ \rho\ m\ M$ give $\mathtt{LE\_Interp}\ \rho\ m'\ M$.
\item $\rho\le\rho'$ preserves interpretations; matching statements for \texttt{Matches},
\texttt{RHS}, \texttt{Const}.
\end{itemize}
\end{theorem}
\begin{proof}
\leanok
\lead{Idea} Induction on the derivation: every rule ends in a $\le$ premise, so shape monotonicity is
transitivity at each constructor; valuation monotonicity pushes through \texttt{Valuation.LE.push} in
the binder cases.
\end{proof}

\begin{lemma}[Shape-level pattern matching is deterministic]
\label{thm:expa-le-interp-matches}
\lean{Lean4Lean.SExpr.LE_Interp.Matches.varN_const_head,Lean4Lean.SExpr.LE_Interp.Matches.arity,Lean4Lean.SExpr.LE_Interp.Matches.head_wf,Lean4Lean.SExpr.LE_Interp.Matches.head_wf_eq,Lean4Lean.SExpr.LE_Interp.Matches.matches_inter,Lean4Lean.SExpr.LE_Interp.Matches.compat_join,Lean4Lean.SExpr.LE_Interp.Matches.unique,Lean4Lean.SExpr.pat_arity,Lean4Lean.SExpr.LE_Interp.Matches.lift}
\leanok
\srcloc{Lean4Lean/Experimental/ShapeLogRel.lean}{3430}
\uses{def:expa-le-interp,inductive:pattern-core}
\thesisref{\texttt{axioms.tex} \S the $\iota$ and $\delta$ rules: \texttt{Pattern}/\texttt{Params.Pat} is this file's abstraction of them.}

\lead{In short} The matcher is well behaved: \hl{same pattern, same result}.
\begin{itemize}
\item The head of a match is the expected constant at the expected arity.
\item Two matches of the same pattern against the same head agree (\texttt{unique}).
\item Two matches of different patterns against compatible arguments force the patterns to intersect
(\texttt{matches\_inter}); two matches of the same pattern against compatible arguments are
compatible and join (\texttt{compat\_join}).
\end{itemize}
\end{lemma}
\begin{proof}
\leanok
\lead{Idea} Induction on the \texttt{Matches} derivation, using \texttt{Params.pat\_wf} (registered
patterns are well formed) and pattern intersection. \texttt{Params.pat\_uniq} first enters at
\srcloc{Lean4Lean/Experimental/ShapeLogRel.lean}{3931}, not here.
\end{proof}

\begin{theorem}[The shapes of a term form a directed set]
\label{thm:expa-le-interp-compat-join}
\lean{Lean4Lean.SExpr.LE_Interp.compat_join,Lean4Lean.SExpr.LE_Interp.compat,Lean4Lean.SExpr.LE_Interp.join,Lean4Lean.SExpr.LE_Interp.join',Lean4Lean.SExpr.LE_Interp.Const.compat_join,Lean4Lean.SExpr.LE_Interp.Const.compat_mismatch,Lean4Lean.SExpr.LE_Interp.RHS.compat_join}
\leanok
\srcloc{Lean4Lean/Experimental/ShapeLogRel.lean}{3936}
\uses{def:expa-le-interp,def:expa-wshape-join}
\lead{In short} Any two shapes approximating the same term \hl{combine into one}.
\begin{itemize}
\item \lead{Given} $\rho'\le\rho$, $\mathtt{LE\_Interp}\ \rho'\ m_1\ M$, $\mathtt{LE\_Interp}\ \rho\
m_2\ M$.
\item \lead{Then} $m_1$, $m_2$ compatible and $\mathtt{LE\_Interp}\ \rho\ (m_1\sqcup m_2)\ M$.
\end{itemize}
\end{theorem}
\begin{proof}
\uses{thm:expa-le-interp-matches,thm:expa-wshape-join-prop,thm:expa-hastype-join,thm:expa-le-interp-mono}
\leanok
\lead{Idea} Four stacked inductions: \texttt{RHS.compat\_join}, \texttt{Const.compat\_mismatch},
\texttt{Const.compat\_join}, then the derivation itself, using \ref{thm:expa-wshape-join-prop} and
\ref{thm:expa-hastype-join}; the hard case is a $\delta$/pattern mismatch between two unfoldings of
the same constant.
\end{proof}

\begin{theorem}[Substitution and weakening for the shape interpretation]
\label{thm:expa-le-interp-subst}
\lean{Lean4Lean.SExpr.LE_Interp.subst,Lean4Lean.SExpr.LE_Interp.inst,Lean4Lean.SExpr.LE_Interp.lift,Lean4Lean.SExpr.LE_Interp.weak,Lean4Lean.SExpr.LE_Interp.weak_iff,Lean4Lean.SExpr.LE_Interp.weak'_iff,Lean4Lean.SExpr.LE_Interp.closed,Lean4Lean.SExpr.LE_Interp.closed_iff}
\leanok
\srcloc{Lean4Lean/Experimental/ShapeLogRel.lean}{4087}
\uses{def:expa-le-interp}
\thesisref{\texttt{typesys.tex} \S properties of substitution (\texttt{thm:subst}): the semantic counterpart.}

\lead{In short} $\mathtt{LE\_Interp}\ \rho\ m\ (M[\sigma])$ \hl{iff a suitable $\rho'$ exists}.
\begin{itemize}
\item $\mathtt{LE\_Interp}\ \rho'\ m\ M$ and $\mathtt{LE\_Interp}\ \rho\ (\rho'\ i)\ (\sigma\ i)$ for
every $i$.
\item Specialised to single instantiation (\texttt{inst}), weakening (\texttt{weak}), closed terms
(\texttt{closed}).
\end{itemize}
\end{theorem}
\begin{proof}
\uses{thm:expa-le-interp-mono,thm:expa-le-interp-compat-join}
\leanok
\lead{Idea} Forward: generalise over a term whose substitution equals the given one, induct, build
$\rho'$, joining sub-valuations by \ref{thm:expa-le-interp-compat-join} at binders. Backward: direct
induction.
\end{proof}

\begin{theorem}[Inversion of the interpretation at Pi and lambda shapes]
\label{thm:expa-le-interp-inv}
\lean{Lean4Lean.SExpr.LE_Interp.forallE_inv,Lean4Lean.SExpr.LE_Interp.forallE_inv',Lean4Lean.SExpr.LE_Interp.lam_inv,Lean4Lean.SExpr.LE_Interp.lam_inv'}
\leanok
\srcloc{Lean4Lean/Experimental/ShapeLogRel.lean}{4241}
\uses{def:expa-le-interp,def:expa-tshape}
\thesisref{\texttt{unique.tex} lines 29--39, clause 2 of definitional inversion; the model-side version, transported to syntax by \ref{thm:expa-lr-adequacy}.}

\lead{In short} If a $\Pi$ or a lambda is interpreted by a matching shape, \hl{so are its parts}.
\begin{itemize}
\item $\forall B.\,F$ interpreted by $\mathtt{forallE}\ b\ f$ gives $b$ interprets $B$ and, for every
$X$ of shape $x$, $f\ @\ x$ interprets $F[X]$.
\item Dually for a lambda against a $\mathtt{lam}$ shape.
\end{itemize}
\end{theorem}
\begin{proof}
\uses{thm:expa-le-interp-subst,thm:expa-le-interp-mono}
\leanok
\lead{Idea} Inversion on the rule; \texttt{TShape.LE.forallE\_decomp} pushes $\le$ into components at a
common level; \texttt{WShapeFun.app\_eq} relates $f\ @\ x$ to the table entry used;
\ref{thm:expa-le-interp-subst} for the codomain.
\end{proof}

\begin{definition}[Fitting valuations and typed interpretation]
\label{def:expa-fits-interptyped}
\lean{Lean4Lean.SExpr.Valuation.Fits,Lean4Lean.SExpr.InterpTyped,Lean4Lean.SExpr.InterpTyped.bot,Lean4Lean.SExpr.InterpTyped.mk,Lean4Lean.SExpr.InterpTyped.out,Lean4Lean.SExpr.InterpTyped.hsort,Lean4Lean.SExpr.InterpTyped.hsort'}
\leanok
\srcloc{Lean4Lean/Experimental/ShapeLogRel.lean}{4323}
\uses{def:expa-le-interp,def:expa-shape-hastype,def:expa-valuation}
\thesisref{\texttt{soundness.tex} Theorem \texttt{thm:sound2}: $\mathtt{InterpTyped}$ is the finitary shadow of $e\in\mathrm{tag}(\alpha)$, \texttt{Fits} that of $\gamma\in\Gamma$.}

\lead{In short} $\mathtt{Fits}$: a valuation assigns each new variable a shape \hl{typed at an
interpretation of its declared type}; $\mathtt{InterpTyped}$ bundles a term, its shape, and that
typing.
\begin{itemize}
\item $\mathtt{Fits}\ \Gamma\ \Delta\ \rho$: $\Delta$ extends $\Gamma$; every assigned shape's type
shape can itself be raised to a type.
\item $\mathtt{InterpTyped}\ \rho\ m\ M\ A$: $\exists m'\ge m, a$ with $m'$ interpreting $M$, $a$
interpreting $A$, $m':a$; \texttt{out} normalises both to a common level.
\end{itemize}
\end{definition}

\section{Soundness of the interpretation}

\begin{theorem}[Typed-interpretation cases for application, lambda and Pi]
\label{thm:expa-le-interp-sound-cases}
\lean{Lean4Lean.SExpr.LE_Interp.sound_bot,Lean4Lean.SExpr.LE_Interp.sound_app,Lean4Lean.SExpr.LE_Interp.sound_lam,Lean4Lean.SExpr.LE_Interp.sound_forallE}
\leanok
\srcloc{Lean4Lean/Experimental/ShapeLogRel.lean}{4361}
\uses{def:expa-fits-interptyped}
\thesisref{\texttt{soundness.tex} Theorem \texttt{thm:sound} (line 216): the corresponding induction cases.}

\lead{In short} The \hl{inductive steps of soundness}, each a hypothetical.
\begin{itemize}
\item $F$ typed at $\forall A.\,B$, $B[X]$ raisable to a type $\Rightarrow$ $F\ X$ typed at $B[X]$
(\texttt{sound\_app}).
\item Analogously \texttt{sound\_lam} (a lambda at $\forall A.\,B$) and \texttt{sound\_forallE} (a
$\Pi$ at $\mathrm{Sort}(\mathrm{imax}\ u\ v)$).
\item \texttt{sound\_bot}: the trivial $\bot$ case.
\end{itemize}
\end{theorem}
\begin{proof}
\uses{thm:expa-hastype-join,thm:expa-le-interp-inv,def:expa-shape-single}
\leanok
\lead{Idea} Discharge $m\le\bot$, invert the interpretation, assemble a witness by joining table
entries; \texttt{sound\_lam}/\texttt{sound\_forallE} induct over the finite table, using singleton
tables (\ref{def:expa-shape-single}) step by step.
\end{proof}

\begin{definition}[Semantic soundness predicates and strong soundness]
\label{def:expa-soundeq-strongsound}
\lean{Lean4Lean.SExpr.SoundEq,Lean4Lean.SExpr.SoundTy,Lean4Lean.SExpr.StrongSound,Lean4Lean.SExpr.StrongSoundCore,Lean4Lean.SExpr.StrongSoundEq}
\leanok
\srcloc{Lean4Lean/Experimental/ShapeLogRel.lean}{4607}
\uses{def:expa-fits-interptyped,def:expa-le-interp}
\thesisref{\texttt{soundness.tex} Theorems \texttt{thm:sound}, \texttt{thm:sound2}: the two bullets are literally \texttt{SoundTy} and \texttt{SoundEq}.}

\lead{In short} $\mathtt{SoundEq}$/$\mathtt{SoundTy}$ state soundness \hl{under every fitting
valuation}; \texttt{StrongSound} strengthens typing to match.
\begin{itemize}
\item $\mathtt{SoundEq}\ \Gamma\ M\ N$: $M$, $N$ have the same shape interpretations.
\item $\mathtt{SoundTy}\ \Gamma\ M\ A$: every shape of $M$ is typed at a shape of $A$.
\item $\mathtt{StrongSound}$/\texttt{StrongSoundCore}: mutually inductive; bundles
$\Gamma\vdash M:A$, \texttt{SoundTy}, and a syntax-directed core derivation at some $A'$ with
$\mathtt{SoundEq}\ \Gamma\ A'\ A$. $\mathtt{StrongSoundEq}$ bundles a defeq with \texttt{SoundEq} and
strong soundness of both sides.
\end{itemize}
\end{definition}

\begin{lemma}[Congruence and equivalence properties of SoundEq]
\label{family:expa-soundeq-lemmas}
\lean{Lean4Lean.SExpr.SoundEq.rfl,Lean4Lean.SExpr.SoundEq.symm,Lean4Lean.SExpr.SoundEq.trans,Lean4Lean.SExpr.SoundEq.inst,Lean4Lean.SExpr.SoundEq.sort,Lean4Lean.SExpr.SoundEq.forallE,Lean4Lean.SExpr.SoundEq.forallE_inv,Lean4Lean.SExpr.SoundTy.defeq_l,Lean4Lean.SExpr.SoundTy.defeq_r,Lean4Lean.SExpr.StrongSoundEq.rfl,Lean4Lean.SExpr.StrongSoundEq.symm,Lean4Lean.SExpr.StrongSoundEq.trans,Lean4Lean.SExpr.StrongSoundEq.mk',Lean4Lean.SExpr.StrongSoundEq.hasType,Lean4Lean.SExpr.StrongSound.defeq,Lean4Lean.SExpr.StrongSound.sound}
\leanok
\srcloc{Lean4Lean/Experimental/ShapeLogRel.lean}{4639}

\lead{In short} \texttt{SoundEq} is an \hl{equivalence relation and a congruence} for sorts and $\Pi$.
\begin{itemize}
\item Preserved by instantiation with a well-typed term; \texttt{forallE\_inv} inverts it at $\Pi$
types.
\item \texttt{SoundTy} transfers along \texttt{SoundEq} on either side; \texttt{StrongSoundEq}
inherits reflexivity, symmetry, transitivity.
\end{itemize}
\end{lemma}
\begin{proof}
\uses{def:expa-soundeq-strongsound,thm:expa-le-interp-subst,thm:expa-le-interp-inv}
\leanok
\lead{Idea} Mostly unfolding, since \texttt{SoundEq} is a biconditional over fitting valuations;
\texttt{inst} uses \ref{thm:expa-le-interp-subst} and \texttt{Fits.cons}; \texttt{forallE\_inv} uses
\ref{thm:expa-le-interp-inv}.
\end{proof}

\begin{theorem}[Semantic unique typing]
\label{thm:expa-strongsound-uniq}
\lean{Lean4Lean.SExpr.StrongSound.uniq,Lean4Lean.SExpr.StrongSound.uniq_of_core}
\leanok
\srcloc{Lean4Lean/Experimental/ShapeLogRel.lean}{4750}
\uses{def:expa-soundeq-strongsound}
\thesisref{\texttt{unique.tex} Theorems \texttt{thm:unique}, \texttt{thm:utype}: proved here semantically, without the $\vdash_n$ stratification.}

\lead{In short} The type of a term is unique \hl{up to sameness of shape interpretation}.
\begin{itemize}
\item \lead{Given} $\mathtt{StrongSound}\ \Gamma\ M\ A$ and $\mathtt{StrongSound}\ \Gamma\ M\ B$.
\item \lead{Then} $\mathtt{SoundEq}\ \Gamma\ A\ B$.
\item \lead{Caveat} Unconditional in itself; supplying \texttt{StrongSound} for an arbitrary
derivation goes through Theorem~\ref{thm:expa-le-interp-strongsound}, which is tainted.
\end{itemize}
\end{theorem}
\begin{proof}
\uses{family:expa-soundeq-lemmas}
\leanok
\lead{Idea} Induction on $M$, reduced by \texttt{uniq\_of\_core} to the syntax-directed core
derivations, matched pairwise: variable by uniqueness of \texttt{Lookup}, application by
\texttt{forallE\_inv}/\texttt{inst}, lambda by \texttt{SoundEq.forallE}, $\Pi$ by
\texttt{SoundEq.sort}.
\end{proof}

\begin{lemma}[Realising application spines against pattern rules]
\label{thm:expa-build-spine}
\lean{Lean4Lean.SExpr.LE_Interp.apps_realize,Lean4Lean.SExpr.LE_Interp.apps_realize_inv,Lean4Lean.SExpr.LE_Interp.RHS.le_applyS,Lean4Lean.SExpr.LE_Interp.RHS.of_applyS,Lean4Lean.SExpr.LE_Interp.Matches.of_matchesS,Lean4Lean.SExpr.LE_Interp.build_spine}
\leanok
\srcloc{Lean4Lean/Experimental/ShapeLogRel.lean}{4764}
\uses{def:expa-le-interp,def:expa-soundeq-strongsound,inductive:pattern-core}
\thesisref{\texttt{axioms.tex} \S the $\iota$ rule: the interpretation-side account of $\iota$/$\delta$ as pattern rewrites.}

\lead{In short} Connects the \hl{syntactic matcher} with its shape-level counterpart.
\begin{itemize}
\item A term matching a registered pattern decomposes into a constant head applied to a shaped spine
(\texttt{build\_spine}, \texttt{Matches.of\_matchesS}).
\item Shape interpretations transfer across a rule's right-hand side, both directions
(\texttt{RHS.le\_applyS}, \texttt{RHS.of\_applyS}).
\end{itemize}
\end{lemma}
\begin{proof}
\uses{thm:expa-le-interp-matches}
\leanok
\lead{Idea} Induction on \texttt{MatchesS}/\texttt{RHS}; \texttt{apps\_realize\_inv} peels one
application at a time off a $\Pi$-telescope, tracking \texttt{StrongSound} typing of the partial
spine.
\end{proof}

\begin{theorem}[Soundness of the shape interpretation]
\label{thm:expa-le-interp-strongsound}
\lean{Lean4Lean.SExpr.LE_Interp.strongSound,Lean4Lean.SExpr.LE_Interp.sound}
\leanok
\srcloc{Lean4Lean/Experimental/ShapeLogRel.lean}{5054}
\uses{def:expa-soundeq-strongsound,inductive:expb-sexpr-isdefeq}
\thesisref{\texttt{soundness.tex} Theorems \texttt{thm:sound}, \texttt{thm:sound2}.}

\lead{In short} Definitional equality \hl{preserves the shape interpretation and its typing}.
\begin{itemize}
\item \lead{Given} $\Gamma\vdash M\equiv N:A$.
\item \lead{Then} $\mathtt{StrongSoundEq}\ \Gamma\ M\ N\ A$; consequently, under every fitting
valuation, $M$ and $N$ have the same shape interpretations and every shape of $M$ is typed at a shape
of $A$ (\texttt{LE\_Interp.sound}).
\end{itemize}
\axfoot{sorryAx via SExpr.IsDefEq.strong (Experimental/SExpr.lean:679) and the SExpr substitution
and weak-head-reduction lemmas; plus the declared axiom Lean4Lean.SExpr.Params.extra\_pat
(Experimental/SExpr.lean:614)}
\end{theorem}
\begin{proof}
\uses{thm:expa-le-interp-sound-cases,thm:expa-build-spine,thm:expa-le-interp-compat-join,thm:expa-hastype-proofirrel,class:expb-sexpr-params}
\leanok
\lead{Idea} Induction on the \emph{strong} defeq derivation (\texttt{replace H := H.strong}): variables by
\texttt{Fits} inversion, $\beta$/$\eta$ by \ref{thm:expa-le-interp-subst}, proof irrelevance by
\ref{thm:expa-hastype-proofirrel}, application/lambda/$\Pi$ by
\ref{thm:expa-le-interp-sound-cases}, the $\delta$/$\iota$ case by \ref{thm:expa-build-spine} and
\texttt{Params.extra\_pat}.
\end{proof}

\section{The shape-indexed logical relation}

\begin{definition}[The LogRel interface]
\label{structure:expa-logrel}
\lean{Lean4Lean.SExpr.LogRelBase,Lean4Lean.SExpr.LogRel,Lean4Lean.SExpr.LogRelBase.TyDefEq.sort}
\leanok
\srcloc{Lean4Lean/Experimental/ShapeLogRel.lean}{5265}
\uses{def:expa-wshape,def:expa-shape-hastype,family:expb-sexpr-operational}
\lead{In short} A record of \hl{nineteen closure obligations} over two predicates on
$\mathtt{WShape}\ n$, packaged so the successor step (\ref{def:expa-lrs}) can be built once,
generically.
\begin{itemize}
\item $\mathtt{DefEq}\ M\ N\ A\ m\ a$: $M$, $N$ equal at type $A$, element-shape $m$, type-shape $a$.
$\mathtt{TyDefEq}\ A\ B\ a$: $A$, $B$ equal types at shape $a$.
\item At a sort: both demand the same weak-head-reduced sort (\texttt{sort\_iff},
\texttt{sort\_iff\_ty}); trivial at $\bot$ when $a:\mathrm{Type}$ (\texttt{bot}); \texttt{toType}
turns a sort-typed \texttt{DefEq} into a \texttt{TyDefEq} at the element shape.
\item Left-reflexivity, symmetry, three transitivities, conversion, four monotonicity fields, closure
of \texttt{TyDefEq} under compatible joins, and invariance under weak-head reduction of both sides.
\end{itemize}
\end{definition}

\begin{definition}[The relation at shape level zero]
\label{def:expa-lr0}
\lean{Lean4Lean.SExpr.LR0,Lean4Lean.SExpr.LR0.TyDefEq,Lean4Lean.SExpr.LR0.DefEq}
\leanok
\srcloc{Lean4Lean/Experimental/ShapeLogRel.lean}{5298}
\uses{structure:expa-logrel,family:expb-sexpr-operational}
\thesisref{\texttt{unique.tex} Theorem \texttt{thm:0dinv} (line 56): the base case, reindexed by shape level.}

\lead{In short} At level $0$ a well-formed shape is $\bot$ or a sort: the relation demands \hl{the
same sort on both sides} -- exactly where sort injectivity enters.
\begin{itemize}
\item $\mathtt{LR0.TyDefEq}\ \Gamma\ M\ N$: $\mathrm{True}$ at $\bot$; at a sort, some
$\mathrm{Sort}\ u$ is a common weak-head reduct of $M$ and $N$.
\item $\mathtt{LR0.DefEq}$ matches on the type-shape: $\mathrm{True}$ at $\bot$,
$\mathtt{LR0.TyDefEq}$ at a sort.
\item All nineteen obligations by \texttt{simp}/\texttt{split}/\texttt{trivial}; the three
transitivities use determinacy of weak-head reduction.
\end{itemize}
\end{definition}

\begin{definition}[Pi-edge and lambda validity at level n+1]
\label{def:expa-lrs-components}
\lean{Lean4Lean.SExpr.LRS.PiDefEq,Lean4Lean.SExpr.LRS.ValTyPi2,Lean4Lean.SExpr.LRS.LamDefEq,Lean4Lean.SExpr.LRS.PiDefEq.left,Lean4Lean.SExpr.LRS.LamDefEq.left}
\leanok
\srcloc{Lean4Lean/Experimental/ShapeLogRel.lean}{5376}
\uses{structure:expa-logrel,def:expa-shape-app}
\lead{In short} Three predicates relative to a predecessor relation $\mathit{IH}:\mathtt{LogRel}\
\Gamma\ n$: \hl{lift $\Pi$/lambda equality one level up}.
\begin{itemize}
\item $\mathtt{PiDefEq}\ \mathit{IH}\ B\ F_1\ F_2\ b\ f$: for related arguments $x,y:B$ at shape
$p:b$, $F_i[x]\equiv F_i[y]$ as types at $f\ @\ p$ ($i=1,2$); plus a single-argument version.
\item $\mathtt{ValTyPi2}$: $M_1,M_2$ weak-head reduce to $\Pi$s with convertible, $b$-typed domains
and codomains convertible under the first domain and satisfying \texttt{PiDefEq}.
\item $\mathtt{LamDefEq}$: the same pattern for terms -- applying $M,N$ to related arguments gives
results related at $(m\ @\ p,\ a_2\ @\ p)$.
\item \texttt{left} lemmas project a two-sided instance to a reflexive one.
\end{itemize}
\end{definition}

\begin{definition}[Type equality at level n+1]
\label{def:expa-lrs-tydefeq}
\lean{Lean4Lean.SExpr.LRS.TyDefEq,Lean4Lean.SExpr.LRS.TyDefEq.left,Lean4Lean.SExpr.LRS.TyDefEq.symm,Lean4Lean.SExpr.LRS.TyDefEq.trans,Lean4Lean.SExpr.LRS.TyDefEq.bot,Lean4Lean.SExpr.LRS.TyDefEq.sort_iff,Lean4Lean.SExpr.LRS.TyDefEq.lam',Lean4Lean.SExpr.LRS.TyDefEq.ctor',Lean4Lean.SExpr.LRS.TyDefEq.forallE_iff}
\leanok
\srcloc{Lean4Lean/Experimental/ShapeLogRel.lean}{5431}
\uses{def:expa-lrs-components,structure:expa-logrel}
\thesisref{\texttt{unique.tex} lines 29--39, clauses 1--2: non-trivial exactly at sort and $\Pi$ shapes.}

\lead{In short} Defined by cases on the element-shape; \hl{trivial except at sort and $\Pi$}, where a
non-type shape carries no type information.
\begin{center}
\begin{tabular}{p{0.42\linewidth}p{0.50\linewidth}}
\toprule
Shape & \texttt{LRS.TyDefEq} \\
\midrule
$\bot$, $\mathtt{lam}$, $\mathtt{ctor}$, $\mathtt{indTy}$ & $\mathrm{True}$ \\
sort & common weak-head normal sort \\
$\mathtt{forallE}\ b\ f$ & $\mathtt{ValTyPi2}\ \mathit{IH}\ M\ N\ b\ f$ \\
\bottomrule
\end{tabular}\par\smallskip
\end{center}
Left-reflexivity, symmetry, transitivity from the corresponding $\mathit{IH}$ fields; transitivity
uses determinacy of weak-head reduction, both transport the codomain equality along the domain
conversion (\texttt{defeqDF\_l}).
\end{definition}

\begin{lemma}[Monotonicity, join and reduction-invariance at level n+1]
\label{family:expa-lrs-mono}
\lean{Lean4Lean.SExpr.LRS.LamDefEq.mono_r_1,Lean4Lean.SExpr.LRS.LamDefEq.mono_r_2,Lean4Lean.SExpr.LRS.LamDefEq.mono_l,Lean4Lean.SExpr.LRS.PiDefEq.mono_r_2,Lean4Lean.SExpr.LRS.PiDefEq.join,Lean4Lean.SExpr.LRS.LamDefEq.whr,Lean4Lean.SExpr.LRS.LamDefEq.symm,Lean4Lean.SExpr.LRS.LamDefEq.trans}
\leanok
\srcloc{Lean4Lean/Experimental/ShapeLogRel.lean}{5403}
\uses{def:expa-lrs-components}
\lead{In short} \texttt{LamDefEq}/\texttt{PiDefEq} \hl{move along the order, join, and survive
weak-head reduction}.
\begin{itemize}
\item Move up (\texttt{mono\_r\_1}) and down (\texttt{mono\_r\_2}) the domain/codomain type-shapes,
given the matching \texttt{HasTypeLam}/\texttt{HasTypePi} side conditions.
\item \texttt{LamDefEq} is antitone in the element-shape (\texttt{mono\_l}).
\item \texttt{PiDefEq} closed under joins of compatible domain/codomain shapes.
\item \texttt{LamDefEq} invariant under weak-head reduction of both terms, symmetric, transitive.
\end{itemize}
\end{lemma}
\begin{proof}
\uses{thm:expa-hastype-mono,thm:expa-hastype-join,def:expa-shape-app}
\leanok
\lead{Idea} Thread the $\mathit{IH}$ obligation through the quantified argument; intermediate shapes from
table inversions and application monotonicity; \texttt{PiDefEq.join} also uses \texttt{Join.mk} and
\ref{thm:expa-hastype-join}.
\end{proof}

\begin{definition}[Term equality at level n+1]
\label{def:expa-lrs-defeq}
\lean{Lean4Lean.SExpr.LRS.DefEq,Lean4Lean.SExpr.LRS.DefEq.lam_forallE,Lean4Lean.SExpr.LRS.DefEq.bot_a,Lean4Lean.SExpr.LRS.DefEq.sort_a,Lean4Lean.SExpr.LRS.DefEq.bot_m,Lean4Lean.SExpr.LRS.DefEq.sort_forallE,Lean4Lean.SExpr.LRS.DefEq.forallE_forallE,Lean4Lean.SExpr.LRS.DefEq.ctor_forallE,Lean4Lean.SExpr.LRS.DefEq.indTy_forallE,Lean4Lean.SExpr.LRS.DefEq.lam_a,Lean4Lean.SExpr.LRS.DefEq.ctor_a,Lean4Lean.SExpr.LRS.DefEq.indTy_a,Lean4Lean.SExpr.LRS.TyDefEq.lam_m,Lean4Lean.SExpr.LRS.TyDefEq.ctor_m,Lean4Lean.SExpr.LRS.TyDefEq.indTy_m}
\leanok
\srcloc{Lean4Lean/Experimental/ShapeLogRel.lean}{5612}
\uses{def:expa-lrs-components,def:expa-lrs-tydefeq}
\lead{In short} Matches first on the \hl{type-shape}, then, at a $\Pi$ type-shape, on the
element-shape.
\begin{center}
\begin{tabular}{p{0.34\linewidth}p{0.59\linewidth}}
\toprule
Type-shape $a$ & \texttt{LRS.DefEq} \\
\midrule
$\bot$, $\mathtt{indTy}$ & $\mathrm{True}$ \\
sort & $\mathtt{LRS.TyDefEq}$ \\
$\mathtt{forallE}$, element $\bot$ & $\mathrm{True}$ \\
$\mathtt{forallE}\ a_1\ a_2$, element $\mathtt{lam}\ m_g$ & $\exists A_1,A_2$ with
$A\to^*_{\mathrm{wh}}\forall A_1.A_2$, both well typed, $\mathtt{TyDefEq}\ A_1\ A_1\ a_1$,
$\mathtt{PiDefEq}\ \mathit{IH}\ A_1\ A_2\ A_2\ a_1\ a_2$,
$\mathtt{LamDefEq}\ \mathit{IH}\ M\ N\ A_1\ A_2\ m_g\ a_1\ a_2$ \\
any other pairing & $\mathrm{False}$ \\
\bottomrule
\end{tabular}\par\smallskip
\end{center}
Fourteen \texttt{@[simp]} lemmas expose these cases.
\end{definition}

\begin{definition}[The successor step LRS]
\label{def:expa-lrs}
\lean{Lean4Lean.SExpr.LRS}
\leanok
\srcloc{Lean4Lean/Experimental/ShapeLogRel.lean}{5657}
\uses{def:expa-lrs-defeq,def:expa-lrs-tydefeq,family:expa-lrs-mono,structure:expa-logrel}
\thesisref{\texttt{unique.tex} Theorem \texttt{thm:1dinv} (line 258): the inductive step, stratified by shape level, needing no Church--Rosser theorem.}

\lead{In short} Turns a $\mathtt{LogRel}\ \Gamma\ n$ into a $\mathtt{LogRel}\ \Gamma\ (n{+}1)$,
discharging \hl{all nineteen obligations}.
\begin{itemize}
\item \texttt{DefEq}/\texttt{TyDefEq} are \ref{def:expa-lrs-defeq}/\ref{def:expa-lrs-tydefeq}; trivial
head forms delegate to their \texttt{@[simp]} lemmas, the sort case to \texttt{LRS.TyDefEq}, the $\Pi$
case to \ref{family:expa-lrs-mono} or the predecessor relation.
\end{itemize}
\end{definition}

\begin{definition}[The relation LR]
\label{def:expa-lr}
\lean{Lean4Lean.SExpr.LR,Lean4Lean.SExpr.LR_zero,Lean4Lean.SExpr.LR_succ}
\leanok
\srcloc{Lean4Lean/Experimental/ShapeLogRel.lean}{5896}
\uses{def:expa-lr0,def:expa-lrs}
\lead{In short} $\mathtt{LR}\ \Gamma:\mathtt{LogRel}\ \Gamma\ n$ by recursion on the shape level:
$\mathtt{LR}_0=\mathtt{LR0}$, $\mathtt{LR}_{n+1}=\mathtt{LRS}(\mathtt{LR}_n)$, with \hl{unfolding
lemmas} \texttt{LR\_zero}, \texttt{LR\_succ}.
\end{definition}

\begin{theorem}[The relation is invariant under shape lifting]
\label{thm:expa-lr-lift}
\lean{Lean4Lean.SExpr.LR.DefEq.lift,Lean4Lean.SExpr.LR.TyDefEq.lift,Lean4Lean.SExpr.LR.lift_succ_aux,Lean4Lean.SExpr.LRS.PiDefEq.lift_aux,Lean4Lean.SExpr.LRS.LamDefEq.lift_aux}
\leanok
\srcloc{Lean4Lean/Experimental/ShapeLogRel.lean}{6046}
\uses{def:expa-lr,def:expa-shape-lift}
\lead{In short} Lifting both shapes to a higher level \hl{does not change the relation}.
\begin{itemize}
\item \lead{Given} $n\le n'$, $m,a:\mathtt{WShape}\ n$ with $m:a$.
\item \lead{Then} $\mathtt{LR.DefEq}\ M\ N\ A\ (\mathtt{lift}_{n'}m)\ (\mathtt{lift}_{n'}a)\iff
\mathtt{LR.DefEq}\ M\ N\ A\ m\ a$; likewise \texttt{TyDefEq} when $m:\mathrm{Type}$.
\end{itemize}
\end{theorem}
\begin{proof}
\uses{thm:expa-hastype-mono,def:expa-lrs-defeq}
\leanok
\lead{Idea} Induction on $n\le n'$; the one-step case is a simultaneous statement about
\texttt{DefEq}/\texttt{TyDefEq}, with the $\Pi$ and lambda cases delegated to private auxiliaries
taking both induction hypotheses explicitly.
\end{proof}

\begin{definition}[Well-formed pairs of substitutions]
\label{def:expa-lr-substwf}
\lean{Lean4Lean.SExpr.LR.Subst1,Lean4Lean.SExpr.LR.SubstWF,Lean4Lean.SExpr.LR.SubstWF.fits,Lean4Lean.SExpr.LR.SubstWF.toSubstEq,Lean4Lean.SExpr.LR.SubstWF.left,Lean4Lean.SExpr.LR.SubstWF.symm}
\leanok
\srcloc{Lean4Lean/Experimental/ShapeLogRel.lean}{6061}
\uses{def:expa-lr,def:expa-fits-interptyped,family:expb-sexpr-operational}
\thesisref{\texttt{typesys.tex} Theorem \texttt{thm:subst}: the context-level prerequisite of the semantic substitution lemma.}

\lead{In short} $\mathtt{SubstWF}$: the \hl{inductive lifting of \texttt{Subst1}} to whole
substitutions, growing a valuation alongside.
\begin{itemize}
\item $\mathtt{Subst1}\ \Gamma_0\ x\ x'\ A_0\ A\ A'\ \rho\ i$: $\Gamma_0\vdash x\equiv x':A$, and for
every type shape $a$ interpreting $A_0$ under $\rho$, $A\equiv A'$ equal at $a$, and every shape
$m:a$ interpreting $\mathtt{bvar}\ i$ relates $x,x'$ at $(m,a)$.
\item $\mathtt{SubstWF}$: built from \texttt{id} and a \texttt{cons} step; projects to
$\mathtt{Valuation.Fits}$ (\texttt{fits}) and to $\mathtt{Ctx.SubstEq}$ (\texttt{toSubstEq}); closed
under the left component and under symmetry.
\end{itemize}
\end{definition}

\section{Adequacy and definitional inversion}

\begin{definition}[The adequacy predicate]
\label{def:expa-lr-adequate}
\lean{Lean4Lean.SExpr.LR.Adequate,Lean4Lean.SExpr.LR.Adequate.bot,Lean4Lean.SExpr.LR.Adequate.fits,Lean4Lean.SExpr.LR.Adequate.refl,Lean4Lean.SExpr.LR.Adequate.left,Lean4Lean.SExpr.LR.Adequate.symm,Lean4Lean.SExpr.LR.Adequate.trans,Lean4Lean.SExpr.LR.Adequate.trans'}
\leanok
\srcloc{Lean4Lean/Experimental/ShapeLogRelAdequacy.lean}{8}
\uses{def:expa-lr-substwf,def:expa-lr,structure:expa-logrel}
\lead{In short} The \hl{closure of the relation over open terms}: related after every well-formed
substitution.
\begin{itemize}
\item $\mathtt{Adequate}\ \Gamma_0\ \Gamma\ \rho\ M\ N\ A\ m\ a$: for every well-formed pair
$\sigma,\sigma'$ from $\Gamma$ to $\Gamma_0$, $M[\sigma],M[\sigma']$ and $N[\sigma],N[\sigma']$ are
related at $(m,a)$ over $A[\sigma]$; and a single-substitution version.
\item Trivial at $m=\bot$ when $a:\mathrm{Type}$; reflexivity, left-projection, symmetry,
\texttt{trans} immediate from \texttt{LogRel} fields; \texttt{trans'} reassembles its two-sided
component from \texttt{W.left} and \texttt{W.symm.left}.
\end{itemize}
\end{definition}

\begin{theorem}[Extending a well-formed substitution pair]
\label{thm:expa-adequate-cons}
\lean{Lean4Lean.SExpr.LR.Adequate.cons}
\leanok
\srcloc{Lean4Lean/Experimental/ShapeLogRelAdequacy.lean}{45}
\uses{def:expa-lr-substwf,def:expa-lr-adequate}
\lead{In short} A related new pair \hl{extends a well-formed substitution pair} by one variable.
\begin{itemize}
\item \lead{Given} $A$ adequate at every shape interpreting it; $\Gamma\vdash A\equiv A':\mathrm{Sort}\
u$; $p:a_1$ with $a_1$ interpreting $A$; $\Gamma_0\vdash x\equiv x':A[\sigma]$ related at $(p,a_1)$.
\item \lead{Then} for every well-formed pair $\sigma,\sigma'$, the extension $\sigma.x,\sigma'.x'$ is
well formed for $A::\Gamma$ with valuation $\rho.\mathrm{push}\ p$.
\end{itemize}
\end{theorem}
\begin{proof}
\uses{thm:expa-lr-lift,thm:expa-le-interp-strongsound,def:expa-tshape,thm:expa-hastype-join}
\leanok
\lead{Idea} The hard obligation is \texttt{Subst1}: lift both shapes to the max of their levels, join
them, and transport the hypothesis along \ref{thm:expa-lr-lift} and the record's monotonicity
fields.
\end{proof}

\begin{lemma}[From term equality at a sort type to type equality]
\label{thm:expa-lr-tovalty}
\lean{Lean4Lean.SExpr.LR.toValTy}
\leanok
\srcloc{Lean4Lean/Experimental/ShapeLogRelAdequacy.lean}{93}
\uses{def:expa-lr,def:expa-tshape}
\lead{In short} \hl{Read off a type equality} from a term equality at $\mathrm{Sort}\ u$.
\begin{itemize}
\item \lead{Given} $M,N$ related at $\mathrm{Sort}\ u$, element-shape $m$, type-shape $a$
interpreting $\mathrm{Sort}\ u$, $m:a$; $b$ a type shape at level $n\le n'$ with $b\le m$ once the
level is forgotten.
\item \lead{Then} $M,N$ are equal types at $b$.
\end{itemize}
\end{lemma}
\begin{proof}
\uses{thm:expa-lr-lift,structure:expa-logrel}
\leanok
\lead{Idea} Move to $m$'s level and the sort shape (\texttt{mono\_r\_1}), apply \texttt{toType}, lower
again (\texttt{mono\_r\_2\_ty}), transport back along \texttt{LR.TyDefEq.lift}.
\end{proof}

\begin{theorem}[Adequacy, the fundamental theorem of the shape logical relation]
\label{thm:expa-lr-adequacy}
\lean{Lean4Lean.SExpr.LR.adequacy}
\leanok
\srcloc{Lean4Lean/Experimental/ShapeLogRelAdequacy.lean}{105}
\uses{def:expa-lr-adequate,def:expa-le-interp,inductive:expb-sexpr-isdefeq}
\thesisref{\texttt{unique.tex} Theorems \texttt{thm:1dinv} and \texttt{thm:ckappa}: this one theorem replaces both, by the model route rather than Church--Rosser.}

\lead{In short} Definitional equality \hl{gives adequacy at any interpreting shape}.
\begin{itemize}
\item \lead{Given} $\Gamma\vdash M\equiv N:A$; shape $m$ interprets $M$, shape $a$ interprets $A$
under $\rho$, $m:a$.
\item \lead{Then} $M,N$ are $\mathtt{Adequate}$ at $(m,a)$: related under all well-formed
substitution pairs.
\end{itemize}
\axfoot{a literal sorry in the const case (ShapeLogRelAdequacy.lean:154); sorryAx via
SExpr.IsDefEq.strong (Experimental/SExpr.lean:679) and the SExpr substitution lemmas; plus the
declared axiom Lean4Lean.SExpr.Params.extra\_pat (Experimental/SExpr.lean:614)}
\end{theorem}
\begin{proof}
\uses{thm:expa-adequate-cons,thm:expa-lr-tovalty,def:expa-lr-substwf,thm:expa-le-interp-strongsound,thm:expa-le-interp-inv,thm:expa-hastype-proofirrel,class:expb-sexpr-params}
\stSorry{} The \texttt{const} case (a bare constant) is a literal \texttt{sorry}; the rest is a
327-line induction on the strong derivation, covering variables by \texttt{SubstWF} induction,
symmetry/transitivity by transporting the interpretation, application/lambda by
\ref{thm:expa-le-interp-inv}, proof irrelevance by \ref{thm:expa-hastype-proofirrel}, and the
$\delta$/$\iota$ case by \texttt{Params.extra\_pat} and the \texttt{whr} field.
\end{proof}

\begin{theorem}[A type equal to a Pi weak-head reduces to a Pi]
\label{thm:expa-forallE-whred-l}
\lean{Lean4Lean.SExpr.forallE_whRed_l}
\leanok
\srcloc{Lean4Lean/Experimental/ShapeLogRelAdequacy.lean}{434}
\uses{thm:expa-lr-adequacy,def:expa-lrs-components}
\thesisref{\texttt{unique.tex} lines 29--39, clause 2 of definitional inversion.}

\lead{In short} \hl{Weak-head reduction is forced} out of a defeq with a $\Pi$ on one side.
\begin{itemize}
\item \lead{Given} $\Gamma\vdash A_0\equiv\forall B_1.F_1:\mathrm{Sort}\ s$.
\item \lead{Then} $\Gamma\vdash A_0\to^*_{\mathrm{wh}}\forall B_0.F_0$ for some $B_0,F_0$, with
$\Gamma\vdash B_0\equiv B_1:\mathrm{Sort}\ u$ and $B_0::\Gamma\vdash F_0\equiv F_1:\mathrm{Sort}\ v$.
\end{itemize}
\axfoot{sorryAx via LR.adequacy (const case, ShapeLogRelAdequacy.lean:154) and
SExpr.IsDefEq.strong; plus the axiom Lean4Lean.SExpr.Params.extra\_pat}
\end{theorem}
\begin{proof}
\uses{thm:expa-le-interp-strongsound,def:expa-shape-hastype}
\leanok
\lead{Idea} Instantiate \ref{thm:expa-lr-adequacy} at the \emph{minimal} non-trivial $\Pi$ shape
$\mathtt{forallE}\ \bot\ \mathtt{ShapeFun.bot}$ -- the least informative shape that still
discriminates the head form -- then read the reduction and the two conversions off the resulting
$\mathtt{ValTyPi2}$.
\end{proof}

\begin{theorem}[Pi injectivity]
\label{thm:expa-forallE-inv}
\lean{Lean4Lean.SExpr.forallE_inv}
\leanok
\srcloc{Lean4Lean/Experimental/ShapeLogRelAdequacy.lean}{448}
\uses{thm:expa-forallE-whred-l}
\thesisref{\texttt{unique.tex} lines 29--39, clause 2, proved there via Church--Rosser inside Theorem \texttt{thm:1dinv}.}

\lead{In short} A $\Pi$ definitionally equal to a $\Pi$ has \hl{equal domains and codomains}.
\begin{itemize}
\item \lead{Given} $\Gamma\vdash(\forall A_0.B_0)\equiv(\forall A_1.B_1):\mathrm{Sort}\ s$.
\item \lead{Then} $\Gamma\vdash A_0\equiv A_1:\mathrm{Sort}\ u$ and
$A_0::\Gamma\vdash B_0\equiv B_1:\mathrm{Sort}\ v$, for some $u,v$.
\end{itemize}
\axfoot{sorryAx via LR.adequacy (const case) and SExpr.IsDefEq.strong; plus the axiom
Lean4Lean.SExpr.Params.extra\_pat}
\end{theorem}
\begin{proof}
\leanok
\lead{Idea} Apply \ref{thm:expa-forallE-whred-l}; a $\Pi$ is already weak-head normal
(\texttt{WHNF.forallE.whRedS}), so the reduction it produces is the identity.
\end{proof}

\begin{theorem}[A sort is never definitionally equal to a Pi]
\label{thm:expa-sort-forallE-inv}
\lean{Lean4Lean.SExpr.sort_forallE_inv}
\leanok
\srcloc{Lean4Lean/Experimental/ShapeLogRelAdequacy.lean}{455}
\uses{thm:expa-forallE-whred-l}
\thesisref{\texttt{unique.tex} lines 29--39, clause 3 of definitional inversion.}

\lead{In short} \hl{No derivation} of $\Gamma\vdash\mathrm{Sort}\ u\equiv\forall
A_1.B_1:\mathrm{Sort}\ s$.
\axfoot{sorryAx via LR.adequacy (const case) and SExpr.IsDefEq.strong; plus the axiom
Lean4Lean.SExpr.Params.extra\_pat}
\end{theorem}
\begin{proof}
\leanok
\lead{Idea} By \ref{thm:expa-forallE-whred-l} the sort would weak-head reduce to a $\Pi$, refuted by
\texttt{WHNF.sort.whRedS}.
\end{proof}

\begin{theorem}[Sort injectivity]
\label{thm:expa-sort-inv}
\lean{Lean4Lean.SExpr.sort_inv}
\leanok
\srcloc{Lean4Lean/Experimental/ShapeLogRelAdequacy.lean}{458}
\uses{thm:expa-lr-adequacy,def:expa-lr0}
\thesisref{\texttt{unique.tex} lines 29--39, clause 1, proved there inside Theorem \texttt{thm:1dinv}.}

\lead{In short} $\mathrm{Sort}\ u\equiv\mathrm{Sort}\ v$ forces \hl{$u=v$} -- as \texttt{SLevel}s, a
semantic quotient, not merely syntactically.
\begin{itemize}
\item \lead{Given} $\Gamma\vdash\mathrm{Sort}\ u\equiv\mathrm{Sort}\ v:V$.
\item \lead{Then} $u=v$.
\end{itemize}
\axfoot{sorryAx via LR.adequacy (const case) and SExpr.IsDefEq.strong; plus the axiom
Lean4Lean.SExpr.Params.extra\_pat}
\end{theorem}
\begin{proof}
\uses{thm:expa-le-interp-strongsound,def:expa-shape-hastype}
\leanok
\lead{Idea} Interpret $\mathrm{Sort}\ u$ by the level-1 shape $\mathtt{sort}(u\neq 0)$, get a type shape
for it via \ref{thm:expa-le-interp-strongsound}, show it must be $\mathtt{sort}\ \mathtt{true}$;
apply \ref{thm:expa-lr-adequacy} at the identity substitution and read \texttt{sort\_iff}: both sides
reduce to the same $\mathrm{Sort}\ w$, and sorts are normal, so $u=w=v$.
\end{proof}

\section{Abandoned branches}

The remaining two modules, \texttt{LogRel.lean} and \texttt{CoinductiveLogRel.lean}, are earlier,
abandoned attempts at the same target. Nothing live imports them, but the
\texttt{Lean4Lean.Experimental} library glob still compiles both, and several of their declaration
names collide with the modules above.

\begin{definition}[Semantic classifiers, logrel-mltt style]
\label{structure:expa-logrel-classifier}
\lean{Lean4Lean.SExpr.Classifier',Lean4Lean.SExpr.Classifier,Lean4Lean.SExpr.Classifier.EqTy,Lean4Lean.SExpr.Classifier.HasTy,Lean4Lean.SExpr.Classifier.DefEq,Lean4Lean.SExpr.Classifier.DefEq.symm,Lean4Lean.SExpr.Classifier.DefEq.left,Lean4Lean.SExpr.Classifier.defEq_self,Lean4Lean.SExpr.Classifier.DefEq.imp}
\leanok
\srcloc{Lean4Lean/Experimental/LogRel.lean}{8}

\lead{In short} A record of \hl{three defaultable predicates} (type equality, inhabitation, term
equality), each conjoined with the matching syntactic judgement so the semantic layer never claims
more than the syntax does.
\begin{itemize}
\item Three obligations: equality symmetric, left-total, reflexive on inhabitants.
\item $\mathtt{Classifier}\ \Gamma\ A\ u$: a phantom-indexed alias; $\mathtt{DefEq.imp}$ transports a
semantic equality along a syntactic map.
\item A same-named \texttt{Classifier'} with a different signature also lives in
\texttt{StepIndexed.lean} (\srcloc{Lean4Lean/Experimental/StepIndexed.lean}{8}).
\end{itemize}
\end{definition}

\begin{definition}[The Pi classifier and normal types]
\label{def:expa-logrel-classifier-forallE}
\lean{Lean4Lean.SExpr.Classifier.forallE,Lean4Lean.SExpr.NormalType,Lean4Lean.SExpr.NormalType.whnf,Lean4Lean.SExpr.NormalType.weak',Lean4Lean.SExpr.NormalType.weak'_inv}
\leanok
\srcloc{Lean4Lean/Experimental/LogRel.lean}{42}
\uses{structure:expa-logrel-classifier}
\lead{In short} Built \hl{Kripke-style}: families of classifiers for the domain and for the
instantiated codomain, indexed by domain inhabitants.
\begin{itemize}
\item A type equals $\forall A.B$ when it weak-head reduces to a $\Pi$ with pointwise related
components; a term inhabits it when it sends related arguments to related results in every extended
context; two terms are equal when they agree on all inhabitants.
\item $\mathtt{NormalType}$: ``is a sort or a $\Pi$'', with its weak-head-normality and weakening
lemmas; declared again, differently, in \texttt{StepIndexed.lean}
(\srcloc{Lean4Lean/Experimental/StepIndexed.lean}{41}).
\end{itemize}
\end{definition}

\begin{definition}[The inductive logical relation over classifiers]
\label{inductive:expa-logrel-ind}
\lean{Lean4Lean.SExpr.LogRel,Lean4Lean.SExpr.LogRel.hasType,Lean4Lean.SExpr.LogRel.eqTy_refl,Lean4Lean.SExpr.LogRel.cast,Lean4Lean.SExpr.LogRel.cast_eqTy,Lean4Lean.SExpr.LogRel.cast_hasTy,Lean4Lean.SExpr.LogRel.cast_defEq}
\leanok
\srcloc{Lean4Lean/Experimental/LogRel.lean}{70}
\uses{def:expa-logrel-classifier-forallE,structure:expa-logrel-classifier}
\lead{In short} A \texttt{Type}-valued family witnessing that a classifier is \hl{canonical for a
type}, by three constructors.
\begin{itemize}
\item \texttt{stuck}: no weak-head reduct of $A$ is a normal type (Lean's theory is not
normalising), classifier trivial.
\item \texttt{sort}: $A$ reduces to a sort. \texttt{forallE}: $A$ reduces to a $\Pi$, with
Kripke-indexed sub-derivations.
\item Derived: every witness implies $A$ is a type, classifies its own type, and transports along a
type equation.
\item \texttt{LogRel} is declared twice more, incompatibly: as the record of
\ref{structure:expa-logrel} (\srcloc{Lean4Lean/Experimental/ShapeLogRel.lean}{5271}) and again in
\texttt{MoreStepIndexed.lean} (\srcloc{Lean4Lean/Experimental/MoreStepIndexed.lean}{371}).
\end{itemize}
\end{definition}

\begin{lemma}[Irrelevance of the logical-relation witness]
\label{thm:expa-lristype-irrel}
\lean{Lean4Lean.SExpr.LRIsType,Lean4Lean.SExpr.LRIsType.rec,Lean4Lean.SExpr.LRIsType.irrel,Lean4Lean.SExpr.LRIsType.irrel',Lean4Lean.SExpr.LREqTy.irrel,Lean4Lean.SExpr.LRDefEq.irrel,Lean4Lean.SExpr.LRIsType.stuck,Lean4Lean.SExpr.LRIsType.sort,Lean4Lean.SExpr.LRIsType.forallE}
\leanok
\srcloc{Lean4Lean/Experimental/LogRel.lean}{157}
\uses{inductive:expa-logrel-ind}
\lead{In short} Writing $\Gamma\Vdash_u A$ for a classifier plus witness, \hl{any two witnesses
agree}.
\begin{itemize}
\item \lead{Given} $J_1:\Gamma\Vdash_u A$, $J_2:\Gamma\Vdash_{u'}A$.
\item \lead{Then} $u=u'$ and $J_1,J_2$ heterogeneously equal; $\Gamma\Vdash_u A$ is a subsingleton,
so $\mathtt{LREqTy}$/$\mathtt{LRDefEq}$ do not depend on the witness.
\end{itemize}
\end{lemma}
\begin{proof}
\uses{structure:expa-logrel-classifier}
\leanok
\lead{Idea} Induction with the hand-written eliminator \texttt{LRIsType.rec} and case analysis on
$J_2$; determinacy of weak-head reduction rules out mismatched constructors.
\end{proof}

\begin{definition}[Weakening of the logical relation]
\label{def:expa-lristype-weak}
\lean{Lean4Lean.SExpr.LRIsType.weak',Lean4Lean.SExpr.LREqTy.weak',Lean4Lean.SExpr.LRDefEq.weak',Lean4Lean.SExpr.LRIsType.cast,Lean4Lean.SExpr.LREqTy.cast,Lean4Lean.SExpr.LRDefEq.cast,Lean4Lean.SExpr.LRHasTy.cast}
\leanok
\srcloc{Lean4Lean/Experimental/LogRel.lean}{205}
\uses{thm:expa-lristype-irrel,inductive:expa-logrel-ind}
\thesisref{\texttt{typesys.tex} Theorem \texttt{thm:weak} (weakening).}

\lead{In short} A witness for $A$ in $\Gamma$ \hl{transports along a context lifting} to a witness
for $A[\rho]$ in $\Delta$.
\begin{itemize}
\item The derived type- and term-equality predicates transport with it.
\end{itemize}
\end{definition}

\begin{lemma}[Symmetry of semantic type equality, unproved]
\label{thm:expa-lreqty-symm-hole}
\lean{Lean4Lean.SExpr.LREqTy.defeq_r,Lean4Lean.SExpr.LREqTy.symm,Lean4Lean.SExpr.LRDefEq.symm,Lean4Lean.SExpr.LREqTy.defeq,Lean4Lean.SExpr.LRDefEq.defeq,Lean4Lean.SExpr.LRDefEq.left,Lean4Lean.SExpr.LRIsType.hasType}
\leanok
\srcloc{Lean4Lean/Experimental/LogRel.lean}{243}
\uses{thm:expa-lristype-irrel}
\lead{In short} Two claims about transporting equality across semantically equal types; \hl{both are
sorries}.
\begin{itemize}
\item \texttt{LREqTy.defeq\_r}: transport a term equality at $A$ to one at any $B$ semantically equal
to $A$.
\item \texttt{LREqTy.symm}: from a witness for $A$ and a semantic equality to $B$, produce a witness
for $B$ equal back to $A$.
\item The surrounding projections and symmetry at a fixed witness are proved.
\end{itemize}
\end{lemma}
\begin{proof}
\uses{def:expa-logrel-classifier-forallE}
\stSorry{} \texttt{defeq\_r} is a bare \texttt{sorry}; all three branches of \texttt{symm}
(\texttt{stuck}, \texttt{sort}, \texttt{forallE}) are \texttt{sorry}. Symmetry of the semantic
type-equality relation is the standard hard obligation of a logrel-mltt development; this is where
the branch stalled.
\end{proof}

\begin{definition}[The semantic context and substitution layer]
\label{def:expa-logrelv}
\lean{Lean4Lean.SExpr.ClassifierV',Lean4Lean.SExpr.ClassifierV,Lean4Lean.SExpr.ClassifierV.IsTy,Lean4Lean.SExpr.LogRelV,Lean4Lean.SExpr.LVIsCtx,Lean4Lean.SExpr.LVIsType,Lean4Lean.SExpr.LVEqSubst,Lean4Lean.SExpr.LVIsCtx.nil,Lean4Lean.SExpr.LVIsCtx.cons,Lean4Lean.SExpr.LVIsCtx.inv,Lean4Lean.SExpr.LVIsCtxInv,Lean4Lean.SExpr.LVIsCtx.cons_inv,Lean4Lean.SExpr.LVEqSubst.left,Lean4Lean.SExpr.LVEqSubst.tail,Lean4Lean.SExpr.LVEqSubst.head,Lean4Lean.SExpr.LVIsType.subst,Lean4Lean.SExpr.LVEqSubst.eqTy,Lean4Lean.SExpr.LVEqTy,Lean4Lean.SExpr.LVHasType,Lean4Lean.SExpr.LVDefEq,Lean4Lean.SExpr.LVDefEq.refl,Lean4Lean.SExpr.LVDefEq.symm}
\leanok
\srcloc{Lean4Lean/Experimental/LogRel.lean}{258}
\uses{thm:expa-lristype-irrel,def:expa-lristype-weak}
\lead{In short} $\mathtt{ClassifierV}\ \Gamma$: which \hl{pairs of substitutions} into $\Gamma$
count as semantically related; built inductively by $\mathtt{LogRelV}$.
\begin{itemize}
\item $\mathtt{LVIsCtx}$ (validated context, with inversion), $\mathtt{LVIsType}$ (semantically valid
type, as a family indexed by related substitutions), $\mathtt{LVEqSubst}$ with head/tail
projections.
\item Open judgements $\mathtt{LVEqTy}$, $\mathtt{LVHasType}$, $\mathtt{LVDefEq}$, with reflexivity
and symmetry.
\end{itemize}
\end{definition}

\begin{lemma}[Fundamental theorem, abandoned]
\label{thm:expa-logrel-fundamental}
\lean{Lean4Lean.SExpr.fundamental}
\leanok
\srcloc{Lean4Lean/Experimental/LogRel.lean}{353}
\uses{def:expa-logrelv}
\thesisref{\texttt{unique.tex} Theorem \texttt{thm:1dinv} and \texttt{soundness.tex} Theorem \texttt{thm:sound} in spirit; nothing here is established.}

\lead{In short} Would say defeq gives, for every validated context, a \hl{semantically equal typed
pair} -- the predecessor of \ref{thm:expa-lr-adequacy}.
\end{lemma}
\begin{proof}
\uses{thm:expa-lreqty-symm-hole}
\stSorry{} Covers only \texttt{bvar} (behind \texttt{stop}), \texttt{symm} (the one rule
discharged), \texttt{trans} (ends in \texttt{sorry}) and \texttt{sort} (behind \texttt{stop}); every
remaining rule falls through a catch-all \texttt{sorry}.
\end{proof}

\begin{definition}[Classifier skeleton for a coinductive relation]
\label{structure:expa-coind-classifier}
\lean{Lean4Lean.VEnv.Classifier',Lean4Lean.VEnv.Classifier}
\leanok
\srcloc{Lean4Lean/Experimental/CoinductiveLogRel.lean}{15}

\lead{In short} A two-field record (\texttt{HasTy} on terms, \texttt{IsTy} on contexts/types/levels)
over \texttt{VExpr}, plus a phantom-indexed alias.
\begin{itemize}
\item \lead{Caveat} The intended content -- a \texttt{coinductive LogRel} block -- sits in a comment
and was never elaborated (Lean 4 has no \texttt{coinductive} command); only this two-field skeleton
compiles.
\end{itemize}
\end{definition}

\section*{Caveats}
\begin{itemize}
\item Theorem~\ref{thm:expa-le-interp-strongsound} onward rests on a \texttt{sorry}
(\texttt{SExpr.IsDefEq.strong}) and \texttt{axiom Params.extra\_pat}; exactly 23 of 818
declarations in \texttt{ShapeLogRel.lean}, all from that point on, are tainted.
\item \texttt{Params.pat\_wf} and \texttt{pat\_uniq}, assumed of registered $\delta$/$\iota$
patterns, have no exhibited instance: no \texttt{Params} value for Lean's actual rules exists in
this repository.
\item Theorem~\ref{thm:expa-lr-adequacy}, the fundamental theorem, has a literal \texttt{sorry} in
its \texttt{const} case; the three injectivity theorems inherit the gap.
\item The abandoned \texttt{LogRel.lean}/\texttt{CoinductiveLogRel.lean} modules hide further
\texttt{sorry}/\texttt{stop} declarations beyond \ref{thm:expa-lreqty-symm-hole} and
\ref{thm:expa-logrel-fundamental}; \texttt{stop} macro-expands to \texttt{repeat sorry}.
\end{itemize}
